\documentclass[]{article}
\usepackage{deluxetable}

\voffset=-0.6truein
\textheight 8.9truein
\textwidth 6.5truein
\oddsidemargin=-0.05truein

\newcommand{\Msun}{M_{\odot}}

\begin{document}

\centerline{\bf\large Manual for FSPS v4.0}
\centerline{July 29, 2026}
\vspace{0.2in}

\noindent{\bf 1. Overview}
\vspace{0.1in}

\noindent
The collection of fortran routines contained in this package allows the user to
compute simple stellar populations (SSPs) for a variety of IMFs and
metallicities, and for a variety of assumptions regarding the morphology of the
horizontal branch, the blue straggler population, the post--AGB phase, and the
location in the HR diagram of the TP-AGB phase.  A variety of simple and
flexible prescriptions for attenuation by dust are also included, as are dust
emission models from diffuse dust and an AGN torus.  From these SSPs the user
may then generate composite stellar populations (CSPs) for a variety of star
formation histories (SFHs) and dust attenuation prescriptions.  Outputs include
the spectra and magnitudes of the SSPs and CSPs at arbitrary redshift. In
addition to these fortran routines a collection of IDL routines are provided
that allow easy manipulation of the output.

The user is strongly encouraged to read Conroy et al. (2009) for an overview of
stellar population synthesis (SPS) and for details regarding this collection of
routines.  This code has been extensively calibrated against a suite of
observational data (for details see Conroy \& Gunn 2010; Johnson et al. 2021;
Park et al. 2025).  The code package is a in essence a highly flexible SPS code
and has therefore come to be called FSPS.

The rest of this manual is organized as follows.  In $\S$1 we briefly discuss
the overall philosophy of the code structure and highlight the main features.
In $\S$2 we present the routines, how they are used, and provide the user with a
basic program that demonstrates the use of the routines.  $\S$3 discusses the
IDL routines provided to read and manipulate the outputs from the fortran
package.  $\S$4 discusses a variety of questions (ok, only one) the user may
have.

For a description of revisions since the initial release of the code, see the
file REVISION\_HISTORY in the \texttt{doc} directory.  Installation instructions
are also provided in that directory.

\vspace{0.1in}
\noindent{\bf 1.1 Downloading the code}
\vspace{0.1in}

\noindent
FSPS is distributed via github: \texttt{github.com/cconroy20/fsps}.  The code
can be checked out with standard github commands.  Emails will be sent to the
mailing list whenever new versions become available (this is one reason why it
is essential for users of this code to be on the mailing list).

\vspace{0.1in}
\noindent{\bf 1.2 Environment variables}
\vspace{0.1in}

\noindent
As mentioned on the webpage, you will need to set an environment variable called
SPS\_HOME that points to the root SPS directory (i.e. the directory that
contains the \texttt{src}, \texttt{OUTPUTS}, etc. directories)

\vspace{0.1in}
\noindent{\bf 1.3 Inclusion of TP-AGB spectra}
\vspace{0.1in}

\noindent
The TP-AGB empirical spectral library does not extend past the rest-frame K-band
and so in previous versions the integrated spectra were not reliable beyond
$\lambda\approx2.4\mu m$.  The TP-AGB spectra are extrapolated blueward with
simpler linear slopes (as advocated in Lancon \& Wood 2002).  The spectra of the
carbon stars are now extrapolated redward with the Aringer et al. (2009)
synthetic carbon star spectral library.  The spectra of the oxygen-rich TP-AGB
stars are now extrapolated with the latest version of the PHOENIX stellar
spectral library (the BT-SETTL library).

\vspace{0.1in}
\noindent{\bf 1.4 Units}
\vspace{0.1in}

\noindent
This code produces two types of files.  The first are spectral files
(\texttt{*.spec}). The spectra are in units of $L_{\odot}/$Hz, i.e. they are
$f_\nu$.  Integrating over frequency generates the bolometric luminosity of the
spectrum.  Wavelengths are in angstroms {\it in vacuum}, and the wavelength
array for the stellar libraries can be found in the files
\texttt{C3K/c3k\_lr\.lambda} and \texttt{MILES/miles\.lambda}.  The full
wavelength array is also now printed in the first line in the \texttt{*.spec}
files.  The spectral resolution of the library is given in km/s dispersion in
the e.g. \texttt{C3K/c3k\_hr.res} files; negative numbers indicate approximate
values since the spectrum comes from a model based on opacity binning. The
second type of file contains magnitudes in a variety of filters.  The magnitude
zero point (i.e., AB or Vega) is set by the variable
\texttt{compute\_vega\_mags} described below.  Both of these output files
contain information on the age, mass, bolometric luminosity, and star formation
rate.  All of these quantities are in the log (base-10), with units of years,
$M_\odot$, $L_\odot$, and $M_\odot$ yr$^{-1}$, respectively.

\vspace{0.1in}
\noindent{\bf 1.5 Redshift Effects}
\vspace{0.1in}

\noindent
The code allows the user to specify a redshift (see Section 3.1.2 below).
Specifying a non-zero redshift affects {\it both} the computed magnitudes and
spectra.  There are two qualitatively different options for redshifting.  See
the parameter \texttt{redshift\_colors} below for more details.  The returned
magnitudes include both the relevant $(1+z)$ factors and the distance modulus.
Also as of v2.5, IGM absorption via Madau (1995) can be turned on and will
attenuate the spectrum (and mags).  This feature is enabled only for non-zero
redshifts.

\vspace{0.1in}
\noindent{\bf 1.6 Filters}
\vspace{0.1in}

\noindent
The current release contains 159 filters.  The filter names can be found in the
file \texttt{FILTER\_LIST} in the \texttt{data} directory.  The actual filter
definitions are in the file \texttt{allfilters.dat} in the same directory.  The
transmission profiles in this file include atmospheric absorption and are in
units of relative response per photon (as opposed for example to relative
response per unit power).  The filters are normalized internally within the
code.  The output magnitudes are listed in the filter order displayed in the
\texttt{FILTER\_LIST} file.  The user is cautioned to use whenever possible the
exact filter transmission curve appropriate for the data being considered.
There is, for example, no such thing as THE $B-$band filter.  Differences of a
few hundredths of a magnitude are common between different definitions of a
given filter such as $B$.

\vspace{0.1in}
\noindent{\bf 1.7 Computation of magnitudes}
\vspace{0.1in}

\noindent
The AB magnitude through a filter $b$, $m_b$, is defined according to the
following formuale:
\begin{eqnarray}
\langle f_\nu \rangle_b &=& \frac{\int R^b_\gamma\,f_\nu\,\,{\rm d\,
  ln}\,\nu}{\int (\nu/\nu_b)^\beta \,R^b_\gamma\,\, {\rm d \,ln}\,\nu} \\
m_b &=& -2.5{\rm log}_{10}(\langle f_\nu \rangle_b)-48.60
\end{eqnarray}
where $f_\nu$ is the spectrum and $R^b_\gamma$ is the relative response per
photon of the filter.  The factor $(\nu/\nu_b)^\beta$ in the demoninator of
Equation 1 may surprise those used to working with optical photometry.  Indeed,
for UV, optical, and near-IR photometry, one usually adopts $\beta=0$ (e.g., for
the GALEX, SDSS, and 2MASS surveys).  However, IR photometry typically assumes a
different calibration.  For example, the {\it IRAS}, {\it Spitzer} IRAC, and
{\it Herschel} PACS and SPIRE magnitudes are frequently quoted assuming
$\beta=1$ while the {\it Spitzer} MIPS filters adopt $\beta=2$ (i.e., a $10^4$ K
blackbody).  The parameter $\nu_b$ is the central wavelength of the filter.

The motivation underlying this convention is that for sources with an intrinsic
spectral slope $\beta$, the bandpass-convolved flux quoted at frequency $\nu_b$
will be precisely the flux at that frequency. For example, observing a $10^4$K
blackbody through the MIPS filters with the above calibration will return a flux
density at the central frequency of the filter that is precisely the true flux
at $\nu_b$. In the end, this is merely a convention, but one that must be
handled carefully when comparing data to models.  See the routine
\texttt{sps\_setup.f90} for implementation.

\vspace{0.1in}
\noindent{\bf 1.8 Interpretation of the $\Delta_L$ and $\Delta_T$
  parameters}
\vspace{0.1in}

\noindent
As described below, when using the Padova isochrones, the user may modify the
bolometric luminosity and effective temperature of the TP-AGB phase by applying
overall shifts in log($L_{\rm bol}$) and log($T_{\rm eff}$) via the parameters
$\Delta_L$ and $\Delta_T$.  These parameters represent shifts with respect to
the best-fit values found in Villaume et al. (2015).

%\vspace{0.4in}

\vspace{0.1in}
\noindent{\bf 1.9 Dust attenuation models}
\vspace{0.1in}

\noindent
The philosophy of the primary dust absorption model closely follows the two
component model of Charlot \& Fall (2000).  In the general case there is dust
associated with young stars (where the defintion of young is defined by the
parameter \texttt{dust\_tesc}) with an optical depth \texttt{dust1} and a
power-law attenuation curve with index \texttt{dust1\_index}.  There is a second
dust component affecting all stars equally and is implemented as a uniform
screen.  This component has an optical depth \texttt{dust2} and an attenuation
curve specified by the parameter \texttt{dust\_type}.  An optional third
component with optical depth \texttt{dust3} that only affects the older stars
can also be included.  Currently supported attenuation models include a
power-law curve with index \texttt{dust\_index}, a Milky-Way-like curve
(following CCM89 but with a variable UV bump strength), a Calzetti et al. (2000)
curve, models from Witt \& Gordon (2000) and Gordon et al. (2003) and the
parameterized models from Kriek \& Conroy (2013) and Reddy et al. (2015).  Note
that one can specify an effective uniform screen (as commonly employed when
implementing the Calzetti curve, although technically one should not think of of
the Calzetti law as representing a uniform screen) by setting
\texttt{dust1}$=0.0$.

\vspace{0.1in}
\noindent{\bf 1.10 Dust emission model}
\vspace{0.1in}

\noindent
FSPS includes two models for dust emission (see the parameter
\texttt{add\_dust\_emission} below) -- the Draine \& Li 2007 model (DL07) and
the THEMIS model (Jones et al. 2013, 2017).  Since there is currently no
publication discussing this aspect of FSPS, a few words regarding the details
and implementation of the dust emission model are in order.

The Draine \& Li (2007; DL07) is a silicate-graphite-PAH grain model. The model
produces dust emission spectra from $1-10^4\mu m$ as a function of the
interstellar radiation field, $U$, expressed in units of the Milky Way radiation
field.  DL07 advocate constructing spectra for entire galaxies by summing up
emission spectra over a range of radiation field strengths, $P(U)dU$
approximated by a delta function at $U_{\rm min}$ and a power-law component from
$U_{\rm min}<U\leq U_{\rm max}$.  As suggested by DL07, we have fixed $U_{\rm
max}=10^6$ and the power-law slope to be $\alpha=2.0$.  We therefore have:
$P(U)dU=(1-\gamma)\delta(U-U_{\rm min})+\gamma A U^{-2}$,
where $(1-\gamma)$ is the fraction of dust mass exposed to starlight intensity
$U_{\rm min}$ and $A$ is a normalization constant.  The DL07 model thus has
three parameters, $U_{\rm min}$, $\gamma$, and $q_{\rm PAH}$, the latter
parameter being the PAH fraction.  These three parameters are contained in the
parameter set and are called \texttt{duste\_umin}, \texttt{duste\_gamma},
\texttt{duste\_qpah}. The user can define each of these parameters in FSPS.

The THEMIS model has a similar setup to DL07 but with different underlying grain
physics, and a wider allowed range in $q_{\rm PAH}$ and $U_{\rm min}$.

\vspace{0.1in}
\noindent{\bf 1.11 Circumstellar dust}
\vspace{0.1in}

\noindent
Circumstellar dust around AGB stars is an option that can be included if the
switch \texttt{add\_agb\_dust\_model} is turned on, which it is by default.  The
circumstellar dust models are a grid of models derived from the radiative
transfer code DUSTY, see Villaume et al. 2015 for details.

\vspace{0.1in}
\noindent{\bf 1.12 AGN dust}
\vspace{0.1in}

\noindent
Dust emission associated with an AGN torus is an option that can be included if
the switch \texttt{add\_agn\_dust} is turned on, which it is by default.  In
order for AGN dust to have an effect one must also set the \texttt{pset}
parameter \texttt{fagn} to a non-zero value; by default this is set to 0.0.  The
AGN dust model is from Nenkova et al. (2008).  Note that the models available
within FSPS are a very small subset of the grids provided by Nenkova et al.;
i.e., aside from the overall amplitude, the only other free parameter is the
$V-$band optical depth of individual clouds.

\vspace{0.1in}
\noindent{\bf 1.13 Nebular emission}
\vspace{0.1in}

\noindent
Nebular continuum and line emission is included within FSPS based on Cloudy
tables provided by Nell Byler.  The details of these grids are presented in
Byler et al. (2017).   This feature is off by default (for reasons of speed),
and the relevant switch is \texttt{add\_neb\_emission}.  One can also
selectively turn off the continuum emission, choose between Cloudy tables that
do or do not include dust, and can vary the gas-phase metallicity and ionization
parameter.

\vspace{0.1in}
\noindent{\bf 1.14 Bursts of star formation}
\vspace{0.1in}

\noindent
At present the user is able to specify a single burst time and burst strength.
In the future it will be possible to specify multiple bursts.  Note that while
the bursts are added into the spectrophotometric outputs, they are {\it not}
included in the output SFR columns.  It is up to the user to add in these
components by hand to the resulting effective SFR.

\vspace{0.1in}
\noindent{\bf 1.15 Output quantities}
\vspace{0.1in}

\noindent
The basic executable routines available in the public distribution of FSPS
outputs two files containing magnitudes and spectra.  FSPS can also output two
other files: color magnitude diagrams in all available filters and spectral
indices in the Lick system with several additional indices.  In order to output
these files the user must modify the syntax calling \texttt{compsp}
(specifically the \texttt{write\_compsp} variable).

\vspace{0.1in}
\noindent{\bf 1.16 BPASS models}
\vspace{0.1in}

\noindent
BPASS stellar population models are included in FSPS.  These are binary stellar
evolution models (v2.2; Eldridge, Stanway et al. 2017) that are made available
as SSP SEDs for a variety of ages, metallicities, and IMFs.  For this reason the
BPASS models are handled differently than the other isochrone models in FSPS.
Specifically, because the models are tabulated directly as SSPs, the user cannot
modify the spectral library nor the adopted IMF.  The BPASS models we have made
available within FSPS are the ``-bin-imf135all\_100'' models, i.e., they assume
a Salpeter IMF with an upper mass cutoff of $100\Msun$.  We have not included
the very lowest metallicity model and have cut off the age grid at 15.8 Gyr.


\vspace{0.2in}
\noindent{\bf 1.17  General program advice}
\vspace{0.1in}

\noindent
The code provided in this package is optimally designed to be integrated into a
larger fortran program.  The code reads all of the libraries into memory and
then utilizes various routines to compute SSPs and CSPs.  The benefit of the
structure of this package is that it allows the user to produce very large
numbers of models relatively quickly (e.g., $\sim10^5$ models in 30 minutes on a
2.66 GHz Intel processor).  It also allows one to easily integrate SPS into
particular science tasks.

For those more interested in generating quick results, rather than using the
more flexible aspects of the code, we also provide a routine that generates
results after prompting the user for input.  This program is called
\texttt{autosps.exe}.  Note that all outputs are placed in the \texttt{OUTPUTS}
directory, included in the tarball, by default.

There is a well-maintained set of Python interfaces to FSPS, originally created
by Dan Foreman-Mackey and now maintained by Ben Johnson, available here:
\texttt{https://github.com/dfm/python-fsps}.


\vspace{0.4in}
\noindent{\bf 2.  Details of Fortran Routines}
\vspace{0.1in}

\noindent{\bf 2.1 The \texttt{sps\_vars.f90} Module: Common
  Variables, Parameters, and Structure Set-up}
\vspace{0.1in}

\noindent
It is recommended that the user read through the \texttt{sps\_vars.f90} module.
This module must be called at the beginning of every program that uses the
routines described below.  It sets up two types of variables: common and
parameters.  Both parameters and common variables can be seen by all routines
and thus do not have to be explicitly passed during a call to a routine.  The
difference between the two is that the parameters are defined once in
\texttt{sps\_vars.f90} and thereafter cannot be changed.  The common variables,
by contrast, can be changed by various routines (see $\S$ 3.1.1).  Be careful
when changing common variables!  The variables in this module are
well-documented.  Most parameters should not be changed.

The \texttt{sps\_vars.f90} module also sets up structures (for those who are not
familiar with fortran structures, they are similar in many ways to structures in
IDL).  There are two main structures that are used extensively in many routines.
One of these is defined to conveniently handle the output of the
\texttt{compsp.f90} routine; it should not be of much concern for most users.
The second structure defines the parameter set.  It is sufficiently important to
warrant a more detailed discussion (see $\S$ 3.1.2).

\vspace{0.1in}
\noindent{\bf 2.1.1  Defining the isochrones, stellar libraries, and
  dust emission model}
\vspace{0.1in}

\noindent
At the top of the \texttt{sps\_vars.f90} module there are lines that start with
\texttt{\#define}.  These lines set switches that tell the code to compile
certain portions of the module depending on which switches are set.  The
switches are binary, with `1'=on and `0'=off. So for example in the standard
release of \texttt{sps\_vars.f90} you will see:

\noindent \\
\noindent
\texttt{\#define MILES 0}\\
\texttt{\#define C3K\_LR 1}\\
\texttt{\#define C3K\_HR 0}\\
\\
\texttt{\#define PADOVA 0}\\
\texttt{\#define MIST 1}\\
\texttt{\#define PARSEC 0}\\
\texttt{\#define BASTI 0}\\
\texttt{\#define GENEVA 0}\\
\texttt{\#define BPASS 0}\\
\\
\texttt{\#define DL07 1}\\
\texttt{\#define THEMIS 0}\\

\noindent
which tells the code to use the C3K low resolution stellar library, the MIST
isochrones, and the DL07 dust emission model. {\it  It is very important that
you `make clean' every time you change one of these switches.}

\noindent\\
\noindent
Additionally, when using the MIST isochrones and the C3K stellar libraries, it
is possible to use $\alpha$-enhanced isochrones and stellar libraries (e.g. Park
et al. 2025). This is done by setting

\noindent\\
\noindent
\texttt{\#define AFE\_FLAG 1}

\noindent\\
\noindent
Users are encouraged to contact the developers regarding the use of this flag,
as it can become memory intensive and requires specialized spectral libraries.
See Park et al. (2025) for more details and contact \texttt{mp2246@cam.ac.uk}
to obtain the $\alpha$-enhanced higher resolution C3K stellar libraries.

\vspace{0.2in}
\noindent{\bf 2.1.2  Common variables \& parameters}
\vspace{0.1in}

\noindent
We briefly describe several of the most important and most used common
variables and parameters.  These parameters are defined in
\texttt{sps\_vars.f90}.

\begin{itemize}
\item \texttt{add\_agb\_dust\_model}
Parameter to turn on/off the AGB circumstellar dust emission model presented in
Villaume et al. (2014).  By default this is function is turned on as of v2.5.
\end{itemize}

\begin{itemize}
\item \texttt{add\_dust\_emission}
Parameter to turn on/off the dust emission model.
{\it Default is on.}
\end{itemize}

\begin{itemize}
\item \texttt{add\_agn\_dust}
Parameter to turn on/off the AGN dust emission model of Nenkova et al. (2008).
\end{itemize}

\begin{itemize}
\item \texttt{add\_igm\_absorption}
Parameter to turn on/off IGM absorption according to Madau (1995). By default
this option is turned off.
\end{itemize}

\begin{itemize}
\item \texttt{add\_neb\_emission}
Parameter to turn on/off the nebular emission model (both continuum and line
emission), based on Cloudy models from Nell Byler.  By default this option is
turned \emph{off}.
\end{itemize}

\begin{itemize}
\item \texttt{add\_neb\_continuum}
Parameter to turn on/off the nebular continuum.  By default this option is
turned on.  \texttt{add\_neb\_emission} must be turned on for this option to
have any effect.
\end{itemize}

\begin{itemize}
\item \texttt{neblineinspec}
Parameter to turn on/off the inclusion of nebular lines in the spectrum.
\texttt{add\_neb\_emission} must be turned on for this option to have any
effect.
{\it Default is on.}
\end{itemize}

\begin{itemize}
\item \texttt{add\_stellar\_remnants}
Parameter to turn on/off the inclusion of stellar remnants in the calculation of
stellar masses.
{\it Default is on.}
\end{itemize}

\begin{itemize}
\item \texttt{compute\_vega\_mags}
A switch that sets the zero points of the magnitude system:
  \begin{itemize}
  \item {\bf 0:}  AB system
  \item {\bf 1:}  Vega system
  \end{itemize}
{\it Default is 0.}
\end{itemize}


\begin{itemize}
\item \texttt{compute\_light\_ages}
A switch that changes the output ages of COMPSP (see below) to bolometric
luminosity weighted ages.
\end{itemize}

\begin{itemize}
\item \texttt{dust\_type}
Common variable defining the attenuation curve for the diffuse dust component:

  \begin{itemize}
  \item {\bf 0:} power--law; see variable \texttt{dust\_index} below.
  \item {\bf 1:} Milky Way extinction law parameterized by Cardelli et al.
    (1989), with variable UV bump strength; see variables \texttt{mwr} and
    \texttt{uvb} below.
  \item {\bf 2:} Calzetti et al. (2000) attenuation curve.  Note that if this
    option is set then the dust attenuation is applied to all starlight equally
    (not split by age), and therefore the only relevant parameter is
    \texttt{dust2} (defined below), which sets the overall normalization (you
    must set \texttt{dust1}$=0.0$ for this to work correctly).
  \item {\bf 3:} allows the user to access a variety of attenuation curve models
    from Witt \& Gordon (2000).  See the parameters \texttt{wgp1},
    \texttt{wgp2}, and \texttt{wgp3} in $\S$3.1.2.  In this case the parameters
    \texttt{dust1} and \texttt{dust2} have no effect because the WG00 models
    specify the full attenuation curve.
  \item {\bf 4: } Kriek \& Conroy (2013) attenuation curve.  In this model the
    slope of the curve, set by the variable \texttt{dust\_index}, is linked to
    the strength of the UV bump.
  \item {\bf 5: } the SMC bar extinction curve from Gordon et al. (2003).
  \item {\bf 6: } Reddy et al. (2015) attenuation curve.
  \end{itemize}

{\it Default is 0.}

\end{itemize}

\begin{itemize}
\item \texttt{imf\_type}
Common variable defining the IMF type:

  \begin{itemize}
  \item {\bf 0:}  Salpeter (1955)
  \item {\bf 1:}  Chabrier (2003)
  \item {\bf 2:}  Kroupa (2001)
  \item {\bf 3:}  van Dokkum (2008)
  \item {\bf 4:}  Dave (2008)
  \item {\bf 5:}  tabulated piece-wise power-law IMF, specified in
    \texttt{imf.dat} file located in the \texttt{data} directory (or specified
    via the parameter \texttt{imf\_filename}; see below).
  \end{itemize}

{\it Default is 2.}
\end{itemize}

\begin{itemize}
\item \texttt{redshift\_colors}
  \begin{itemize}
  \item {\bf 0:}  Magnitudes are computed at a fixed redshift specified in the
    parameter set (see below)
  \item {\bf 1:}  Magnitudes are computed at a redshift that corresponds to the
    age of the output SSP/CSP (assuming a redshift--age relation appropriate for
    a WMAP7 cosmology).  This switch is useful if the user wants to compute the
    evolution in {\it observed} colors of a SSP/CSP.
  \end{itemize}
\end{itemize}

\begin{itemize}
\item \texttt{smooth\_velocity}
Switch to smooth the spectrum in velocity space.  If off, the spectrum is
smoothed in wavelength space.  The degree of smoothing is determined by the
variable \texttt{sigma\_smooth} in the parameter set (see below).
\end{itemize}

\begin{itemize}
\item \texttt{smooth\_lsf}
Parameter to smooth the SSPs by a tabulated instrumental LSF provided in
data/lsf.dat.  The assumed units are km/s (sigma, not FWHM). Turned off by
default.
\end{itemize}

\begin{itemize}
\item \texttt{tpagb\_norm\_type}
Parameter to choose the normalization of the TP-AGB stars (only applies to the
Padova isochrones; default is 2):

  \begin{itemize}
  \item {\bf 0:}  Original normalization of the Padova isochrones
  \item {\bf 1:}  Normalization from Conroy \& Gunn (2010)
  \item {\bf 2:}  Normalization from Villaume et al. (2015)
  \end{itemize}
\end{itemize}

\begin{itemize}
\item \texttt{time\_res\_incr}
Parameter setting the factor by which the resolution of the default isochrone
age array is increased.  Default value is 1, which is sufficient for most
purposes.
\end{itemize}

\begin{itemize}
\item \texttt{vactoair\_flag}
Parameter to write wavelengths in air rather than vacuum if set to 1.
The default is to write wavelengths in vacuum.
\end{itemize}

\begin{itemize}
\item \texttt{verbose}
Parameter that controls output to screen.  If set to 1 a lot of output will be
printed to screen, if set to 0 the programs will be silent.
\end{itemize}

\clearpage

\vspace{0.1in}
\noindent{\bf 2.1.3  The Parameter Set}
\vspace{0.1in}

\noindent
The parameter set is a structure defined in the \texttt{sps\_vars.f90} module.
It must be defined at the beginning of every program for the various routines to
properly work.  This structure acts as the primary interface between what the
user would like to compute and the various subroutines that actually do the
work.  Simply defining the structure at the beginning of the program will set
all structure elements to their default values (see the example program
\texttt{simple.f90}). The elements of the structure, and the default values, are
described below.  The parameters are organized by topic, e.g., star formation
history, IMF, dust, and stellar evolution parameters.

\vspace{0.3cm}

\noindent
{\bf Basic Parameters}

\begin{itemize}

\item{\texttt{zred}} Redshift.
  If this value is non-zero and if \texttt{redshift\_colors}$=0$, the magnitudes
  will be computed for the spectrum placed at redshift \texttt{zred}.
  {\it Default value is 0.0}.
\item{\texttt{zmet}} Metallicity index.
  The metallicity is specified as an integer ranging between 1 and \texttt{NZ} for the
  isochrones.  A lookup table for the actual metallicities corresponding to the
  integer values is provided at the end of this manual. Note that
  $Z_\odot=0.0190$.
  {\it Default value is 1}.
\item{\texttt{afeindx}} Alpha-enhancement index.
  The $\alpha$ enhancement is specified as an integer ranging between 1 and
  \texttt{N\_AFE} for the isochrones.  This can only be changed if the
  isochrones and stellar libraries are $\alpha$-enhanced (see $\S$2.1.1).
  {\it Default value is 1}.
\item{\texttt{imf1}} Logarithmic slope of the IMF over the range
  $0.08<M<0.5\,\Msun$.  Only used if \texttt{imf\_type}$=2$.
  {\it Default value is 1.3}.
\item{\texttt{imf2}} Logarithmic slope of the IMF over the range
  $0.5<M<1.0\,\Msun$.  Only used if \texttt{imf\_type}$=2$.
  {\it Default value is 2.3}.
\item{\texttt{imf3}} Logarithmic slope of the IMF over the range
  $1.0<M<100\,\Msun$.  Only used if \texttt{imf\_type}$=2$.
  {\it Default value is 2.3}.
\item{\texttt{vdmc}} IMF parameter defined in van Dokkum 2008.  Only
  used if \texttt{imf\_type}$=3$.
  {\it Default value is 0.08}.
\item{\texttt{mdave}} IMF parameter defined in Dave 2008.  Only
  used if \texttt{imf\_type}$=4$.  {\it Default value is 0.5}.
\item{\texttt{evtype}} Only include isochrone points at a particular
  evolutionary phase specified in the isochrone table (see
  \texttt{data/ev\_phases.tex} for phase options).  Only available with for
  isochrones with phase information.  All phases used when set to -1.
  {\it Default value is -1}.
\item{\texttt{masscut}} Only include masses below {\texttt{masscut} in the
  synthesis.
  {\it Default value is 150.0}.
\item{\texttt{imf\_filename}} Filename of the tabulated IMF file located in the
  \texttt{data} directory.  If unset, the default \texttt{imf.dat} is assumed.
  Only applies if \texttt{imf\_type}$=5$.
\item{\texttt{mag\_compute}} Integer array of length \texttt{nbands}. Allows the
  user to choose which magnitudes are computed from the filter list.  1=yes,
  0=no.  All magnitudes are still printed to the mag file, but if the particular
  band has a value of 0 in this array, the output value will be 99.
  {\it Default values are 1 for all bands}.
\item{\texttt{sigma\_smooth} Broadening of the spectrum, in units of km s$^{-1}$
  (if \texttt{smooth\_velocity}$=1$) or \AA, (if \texttt{smooth\_velocity}$=0$).
  This specifies the width of the Gaussian in terms of $\sigma$, {\it not} FWHM.
  {\it Default value is 0.0}.}
\item{\texttt{min\_wave\_smooth} Minimum wavelength, in \AA, that will be
  smoothed if \texttt{sigma\_smooth}$>0$.
  {\it Default value is 1000\,\AA}.}
\item{\texttt{max\_wave\_smooth} Maximum wavelength, in \AA, that will be
  smoothed if \texttt{sigma\_smooth}$>0$.
  {\it Default value is 10,000\,\AA}.}
\item{\texttt{gas\_logu} Value of the gas ionization parameter, $U$; relevant
  only for the nebular emission model.
  {\it Default value is -2.0}.}
\item{\texttt{gas\_logz} Value of the gas metallicity; relevant only for the
  nebular emission model.
  {\it Default value is 0.0}.}
\item{\texttt{igm\_factor}} Factor by which to multiply the default IGM
  absorption optical depth.
    {\it Default value is 1.0}.}
\end{itemize}

\vspace{0.5cm}

\noindent
{\bf SFH Parameters}
\begin{itemize}
\item \texttt{sfh} Defines the type of star formation history, normalized such
  that one solar mass of stars is formed over the full SFH (except for tabular).
  {\it Default value is 0}.
  \begin{itemize}
  \item {\bf 0:} SSP
  \item {\bf 1:} A six parameter SFH (tau model plus a constant component and a
    burst), with parameters \texttt{tau}, \texttt{const}, \texttt{sf\_start},
    \texttt{sf\_trunc}, \texttt{tburst}, and \texttt{fburst} (see below).
  \item {\bf 2:} Tabulated SFH defined in a file called \texttt{sfh.dat} that
    must reside in the \texttt{data} directory (or a file specified via the
    parameter \texttt{sfh\_filename}; see below). The file must contain three
    columns.  The first column is time since the Big Bang in Gyr, the second is
    the SFR in units of solar masses per year, the third is the absolute
    metallicity.  An example is provided in the \texttt{data} directory.  The
    time grid in this file can be arbitrary (so long as the units are correct),
    but it is up to the user to ensure that the tabulated sfh is well-sampled so
    that the outputs are stable. Obviously, highly oscillatory data require
    dense sampling.
  \item {\bf 3:} Reserved for special use of the tabulated SFH option. Allows
    the user to read in the tabulated SFHs directly into the necessary arrays,
    bypassing the need to create \texttt{sfh.dat} files.  Email
    cconroy@cfa.harvard.edu if you would like further instructions for how to
    use this option.
  \item {\bf 4:} This is the same as option 1 except that the tau-model
    component is replaced with a delayed tau model of the form $t\,e^{-t/\tau}$.
  \item {\bf 5:} Delayed tau model with a transition at a time
    \texttt{sf\_trunc} to a linearly decreasing SFH with the slope specified by
    \texttt{sf\_slope}. See Simha et al. 2014 for details.
  \end{itemize}
\item{\texttt{tau}} Defines e-folding time for the SFH, in Gyr.  Only used if
  \texttt{sfh}$=1$ or 4.  The range is $0.1<\tau<10^2$.
  {\it Default value is 1.0}.
\item{\texttt{const}} Defines the constant component of the SFH.  This quantity
  is defined as the fraction of mass formed in a constant mode of SF; the range
  is therefore $0\leq C\leq1$.  Only used if \texttt{sfh}$=1$ or 4.
  {\it Default value is 0.0}.
\item{\texttt{sf\_start}} Start time of the SFH, in Gyr.
  {\it Default value is 0.0}.
\item{\texttt{sf\_trunc}} Truncation time of the SFH, in Gyr.  If set to 0.0,
  there is no trunction.
  {\it Default value is 0.0}.
\item{\texttt{tage}} If set to a non-zero value, the \texttt{compsp} routine
  will compute the spectra and magnitudes only at this age, and will therefore
  only output one age result.  The units are Gyr. (The default is for
  \texttt{compsp} to compute and return results from $t\approx0$ to the maximum
  age in the isochrones).
  {\it Default value is 0.0.}
\item{\texttt{fburst}} Defines the fraction of mass formed in an instantaneous
  burst of star formation.  Only used if \texttt{sfh}$=1$ or 4.
  {\it Default value is 0.0}.
\item{\texttt{tburst}} Defines the age of the Universe when the burst occurs, in
  Gyr.  If \texttt{tburst}$>$\texttt{tage} then there is no burst. Only used if
  \texttt{sfh}$=1$ or 4.  {\it Default value is 11.0}.
\item{\texttt{sf\_slope}} For \texttt{sfh}$=5$, this is the slope of the SFR
  after time \texttt{sf\_trunc}.
  {\it Default value is 0.0}.
\item{\texttt{sfh\_filename}} Filename of the tabulated SFH file located in the
  \texttt{data} directory.  If unset, the default \texttt{sfh.dat} is assumed.
  Only applies if \texttt{sfh}$=2$.
\end{itemize}

\vspace{0.5cm}

\noindent
{\bf Dust Parameters}
\begin{itemize}
\item{\texttt{dust\_tesc}} Stars younger than \texttt{dust\_tesc} are attenuated
  by both \texttt{dust1} and \texttt{dust2}, while stars older are attenuated by
  \texttt{dust2} only.  Units are log(yrs).
  {\it Default value is 7.0}.
\item{\texttt{dust1}} Dust parameter describing the attenuation of young stellar
  light, i.e. where $t\leq$\texttt{dust\_tesc} (for details, see Conroy et al.
  2009a).  Specifically, it is the opacity at 5500\AA.
  {\it Default value is 0.0}.
\item{\texttt{dust2}} Dust parameter describing the attenuation of all stellar
  light.  Specifically, it is the opacity at 5500\AA.
  {\it Default value is 0.0}.
\item{\texttt{dust3}} Optical depth of dust affectting only old stellar light,
  i.e. where $t>$\texttt{dust\_tesc}
  {\it Default value is 0.0}.
\item{\texttt{frac\_nodust}} Fraction of starlight that is not attenuated by the
  diffuse dust component (i.e. that is not affected by \texttt{dust2}).
  {\it Default value is 0.0}.
\item{\texttt{frac\_obrun}} Fraction of starlight that is not attenuated by the
  birth cloud dust component (i.e. that is not affected by \texttt{dust1}).
  {\it Default value is 0.0}.
\item{\texttt{dust\_index}} Power--law index of the diffuse dust attenuation
  curve.  Only used when \texttt{dust\_type}=0 or 4.
  {\it Default value is -0.7}.
\item{\texttt{dust1\_index}} Power--law index of the birth cloud dust
  attenuation curve.  Used for all dust types.
  {\it Default value is -1.0}.
\item{\texttt{mwr}} The ratio of total to selective absorption which
  characterizes the MW extinction curve: $R\equiv A_V/E(B-V)$.  Only used when
  \texttt{dust\_type}=1.
  {\it Default value is 3.1}.
\item{\texttt{uvb}} Parameter characterizing the strength of the 2175\AA\,
  extinction feature with respect to the standard Cardelli et al.  determination
  for the MW.  Only used when \texttt{dust\_type}=1.
  {\it Default value is 1.0}.
\item{\texttt{wgp1}} Integer specifying the optical depth in the Witt \& Gordon
  2000 (WG00) models.  Values range from $1-18$, corresponding to optical depths
  of\\
  $0.25,0.50,0.75,1.00,1.50,2.00,2.50,3.00,3.50,4.00,4.50, 5.00,5.50,6.00,7.00,8.00,9.00,10.0$.
  Note that these optical depths are defined differently from the optical depths
  defined by the parameters \texttt{dust1} and \texttt{dust2}.  See WG00 for
  details.
\item{\texttt{wgp2}} Integer specifying the type of large--scale geometry and
  extinction curve.  Values range from $1-6$, corresponding to MW+dusty,
  MW+shell, MW+cloudy, SMC+dusty, SMC+shell, SMC+cloudy.  MW= Milky Way
  extinction; SMC= Small Magellanic Cloud extinction.  Dusty, shell, and cloudy
  specify the geometry and are described in WG00.
\item{\texttt{wgp3}} Integer specifying the local geometry for the WG00 dust
  models.  A value of 1 corresponds to a homogeneous distribution, and a value
  of 2 corresponds to a clumpy distribution. See WG00 for details.
\item{\texttt{duste\_gamma}} Parameter of the dust emission model.  Specifies
  the relative contribution of dust heated at a radiation field strength of
  $U_{\rm min}$ and dust heated at $U_{\rm min}<U\leq U_{\rm max}$.  Allowable
  range is $0.0-1.0$.
  {\it Default value is 0.01.}
\item{\texttt{duste\_umin}} Parameter of the dust emission model.  Specifies the
  minimum radiation field strength in units of the MW value.  Valid range is
  between 0.1 and 25.0.
  {\it Default value is 1.0.}
\item{\texttt{duste\_qpah}} Parameter of the dust emission model.  Specifies the
  grain size distribution through the fraction of grain mass in PAHs. This
  parameter has units of \% and a valid range of $0.47-4.58$.
  {\it Default value is 3.5.}
\item{\texttt{fagn}} This parameter is the luminosity of the AGN expressed as a
  fraction of the stellar bolometric luminosity.
  {\it Default value is 0.0.}
\item{\texttt{agn\_tau}} $V-$band optical depth of individual clouds in the AGN
  model of Nenkova et al. (2008).
  {\it Default value is 10.0.}
\end{itemize}

\noindent
{\bf Isochrone Parameters}
\begin{itemize}
\item{\texttt{redgb}} Modify weight given to the RGB. Only availble for
  isochrone sets that specify this phase (e.g., BaSTI).
  {\it Default value is 1.0}.
\item{\texttt{agb}} Modify weight given to the AGB.  Only available for
  isochrone sets that specify this phase (e.g., BaSTI).
  {\it Default value is 1.0}.
\item{\texttt{dell}} Shift in ${\rm log}(L_{\rm bol})$ of the TP-AGB isochrones.
  Note that the meaning of this parameter and the one below has changed to
  reflect the updated calibrations presented in Conroy \& Gunn (2009).  That is,
  these parameters now refer to a modification about the calibrations presented
  in that paper.
  {\it Default value is 0.0}.
\item{\texttt{delt}} Shift in ${\rm log}(T_{\rm eff})$ of the TP-AGB isochrones
  for Padova models.
  {\it Default value is 0.0}.
\item{\texttt{agb\_dust}} Scales the AGB circumstellar dust emission model
  presented in Villaume et al. (2014).
  {\it Default value is 1.0}.
\item{\texttt{fcstar}} Currently has no effect.  Previously, the fraction of
  stars that the Padova isochrones identify as Carbon stars that FSPS assigns to
  a Carbon star spectrum.  Set this to 0.0 if for example the users wishes to
  turn all Carbon stars into regular M-type stars.  Valid range is $0.0-1.0$.
  {\it Default value is 1.0}.
\item{\texttt{sbss}} Specific frequency of blue straggler stars.  See Conroy et
  al. (2009) for details and a plausible range.
  {\it Default value is 0.0}.
\item{\texttt{fbhb}} Fraction of horizontal branch stars that are blue.  The
  blue HB stars are uniformly spread in ${\rm log}(T_{\rm eff})$ to $10^4$ K.
  See Conroy et al. 2009a for details and a plausible range.
  {\it Default value is 0.0}.
\item{\texttt{pagb}} Weight given to the post--AGB phase.  A value of 0.0 turns
  off post--AGB stars; a value of 1.0 implies that the Vassiliadis \& Wood 1994
  tracks are implemented as--is.
  {\it Default value is 1.0}.

\end{itemize}

\vspace{0.1in}
\noindent{\bf 2.2  Description of Fortran Routines}
\vspace{0.1in}

\noindent
We now discuss the purpose and syntax of the routines in this package.
$\S$3.2.1 discusses the example routines provided in FSPS that demonstrate how
many of these routines are used.  Routines included in the \texttt{src}
directory that are {\it not} discussed below are meant to only be accessed
through other routines.

%\bigskip

\begin{itemize}
\item \texttt{COMPSP(write\_compsp,nzin,outfile,mass\_ssp,lbol\_ssp,spec\_ssp,pset,ocompsp)}
\end{itemize}

This routine takes as input \texttt{mass\_ssp}, \texttt{lbol\_ssp}, and
\texttt{spec\_ssp}, which are the outputs of the routine \texttt{ssp\_gen.f90}.
The user must also provide the parameter set \texttt{pset} and filename for
output in \texttt{outfile}, if output is desired (a blank string may be
specified if no output is desired). There are six possible outputs, specified by
\texttt{write\_compsp}: No output (0), output magnitudes (1), output spectra
(2), output magnitudes and spectra (3), output spectral indices (4), output
color-magnitude diagrams (5).  The outputs are written to files with the output
filename with ``.mags'', ``.spec'', ``.indx'', or ``.cmd'' appended.

A variety of output from this routine is saved in the \texttt{ocompsp} variable,
which must be a structure as defined in \texttt{sps\_vars.f90}.  Again, see the
example routine below.

The \texttt{nzin} parameter sets the number of metallicity points passed.  For
standard, single metallicity calculations, set this value to one and simply pass
the outputs from \texttt{ssp\_gen.f90}. However, if one wants to compute spectra
for an evolving metallicity (i.e. when computing a tabulated SFH, the
metallicity history may be specified), one needs to set this parameter to the
number of metallicity elements available (in the default release this is 22).
One then must take care to pass the metallicity--dependent \texttt{ssp\_gen.f90}
outputs.  Currently the code only allows the specification of a metallicity
history when passing tabulated SFH.


\begin{itemize}
\item \texttt{GETMAGS(zred,spec,mags)}
\end{itemize}

The input is the redshift \texttt{zred} and spectrum \texttt{spec}. The output
is an array of magnitudes \texttt{mags} for the redshifted spectrum.

\begin{itemize}
\item \texttt{GETINDX(lambda,spec,indices)}
\end{itemize}

This routine computes the spectral indices for the input spectrum \texttt{spec}
with corresponding wavelength array \texttt{lambda}, returning the indices in an
array \texttt{indices}.  The \texttt{indices} array must be defined with an
array length specified by the variable \texttt{nindx}, which specifies the
number of indices.  The indices are defined in the file allindices.dat in the
\texttt{data} directory.

\begin{itemize}
\item \texttt{PZ\_CONVOL(pset,zave,spec\_pz,lbol\_pz,mass\_pz)}
\end{itemize}

This routine convolves the full array of metallicity-dependent SSPs with a
closed-box metallicity distribution function (MDF).  The yield
(\texttt{pset\%pmetals}) is the only input.  Outputs include the average
metallicity (\texttt{zave}) and the spectra, lbol, and mass integrated over the
MDF.  The full metallicity-dependent SSPs must have been previously set up.  An
example of how to use this routine is provided in \texttt{lesssimple.f90}.


\begin{itemize}
\item \texttt{SFHSTAT(pos,model,ssfr6,ssfr7,ssfr8,ave\_age)}
\end{itemize}

This routine returns basic statistics for a given star formation history.  The
inputs are the parameter set (\texttt{pos}) and a single element output from the
compsp routine (\texttt{model}).  The outputs are the specific SFR (SSFR),
averaged over $10^6$, $10^7$, and $10^8$ yrs (\texttt{ssfr6, ssfr7, ssfr8}), and
the mass-weighted average stellar age (\texttt{ave\_age}).  An example of how to
use this routine is provided in \texttt{simple.f90}.


\begin{itemize}
\item \texttt{SPS\_SETUP(zin)}
\end{itemize}

This routine must be called at least once before running any routines. It reads
in all of the isochrones and spectral libraries and stores them in a common
block.  If the user requires only one metallicity, then that metallicity can be
specified as \texttt{zin}.  The metallicity must be specified as an integer
corresponding to the look-up table at the end of this manual.  If the user
wishes to read in all metallicities, then \texttt{zin} should be set to -1.

\begin{itemize}
\item \texttt{SSP\_GEN(pset,mass\_ssp,lbol\_ssp,spec\_ssp)}
\end{itemize}

This routine takes as input the parameter set, \texttt{pset}, and outputs the
time-dependent mass, \texttt{mass\_ssp}, bolometric luminosity,
\texttt{lbol\_ssp}, and spectrum, \texttt{spec\_ssp}, of an SSP defined by the
parameter set.  Each of these input variables must be properly defined at the
beginning of the main routine.  See the example routines for guidance.


\begin{itemize}
\item \texttt{SMOOTHSPEC(lambda,spec,sigma,minl,maxl)}
\end{itemize}

This routine broadens the input spectrum \texttt{spec} with corresponding
wavelength array \texttt{lambda} by a velocity dispersion \texttt{sigma}
measured in km/s.  Note that the velocity broadening is only approximate in that
we are ignoring the variation in d$\lambda$ with $\lambda$ within each
integration step.

\begin{itemize}
\item \texttt{WRITE\_ISOCHRONE(outfile,zz)}
\end{itemize}

This routine writes the isochrones for all ages at a single metallicity to a
file with name \texttt{outfile}.  The metallicity of the isochrone is specified
via the variable \texttt{zz} in the internal integer metallicity units (see the
Table at the end of this manual).  At each age the isochrone is written for all
of the magnitudes specified in the \texttt{FILTER\_LIST} file.

\vspace{0.1in}
\noindent{\bf 2.2.1. Example Routines}
\vspace{0.1in}

The code package contains several routines that highlight many of the features
of FSPS.  The routine \texttt{simple.f90} demonstrates the basic syntax required
to generate simple models.  In addition, there is a less simple routine located
in the \texttt{src} directory, called \texttt{lesssimple.f90}, that highlights
some more advanced features of the code.


\vspace{0.1in}
\noindent{\bf 3.  Description of IDL Routines}
\vspace{0.1in}

\begin{itemize}
\item \texttt{res = read\_indx(file)}
\end{itemize}

This function takes as input an index file and reads it into a simple IDL
structure.  The index file must be created by the user (unlike the .spec and
.mags files, which are automatically created in the compsp routine).  The format
is the age followed by all of the indices defined in the file
\texttt{INDEX\_LIST}.

\begin{itemize}
\item \texttt{res = read\_mags(file)}
\end{itemize}

This function takes as input the magnitude file (*.mags) produced by
\texttt{compsp.f90}.  The output is an IDL structure with elements including the
some of the magnitudes listed in \texttt{FILTER\_LIST} and computed by
\texttt{compsp.f90}.  Note that not all magnitudes computed and listed in the
*.mags file are contained in the resulting structure.  This routine can be
trivially modified to include other/all of the magnitudes computed.

\begin{itemize}
\item \texttt{res = read\_spec(file)}
\end{itemize}

This function takes as input the spectra file (*.spec) produced by
\texttt{compsp.f90}.  The output is an IDL structure with elements including the
time-dependent spectrum computed by \texttt{compsp.f90}.

\begin{itemize}
\item \texttt{res = read\_cmd(file)}
\end{itemize}

This function takes as input the CMD file (*.cmd) produced by
\texttt{compsp.f90}.  The output is an IDL structure with elements including the
time-dependent spectrum computed by \texttt{compsp.f90}.

\begin{itemize}
\item \texttt{res = read\_indx(file)}
\end{itemize}

This function takes as input the index file (*.indx) produced by
\texttt{compsp.f90}.  The output is an IDL structure with elements including the
time-dependent spectrum computed by \texttt{compsp.f90}.

\begin{itemize}
\item \texttt{res = read\_fsps(file)}
\end{itemize}

This function takes as input {\it any} file produced by \texttt{compsp.f90}.
The function determines which type of file was passed based on the filename
(*.spec, *.mags, etc).


\vspace{0.4in}
\noindent{\bf 4.  How do I....}
\vspace{0.1in}

\noindent{\bf 4.1  add additional filters?}
\vspace{0.1in}

\noindent
Adding additional filters is straightforward.  There are three things the user
must do: 1) modify the \texttt{nbands} parameter in the \texttt{sps\_vars.f90}
routine; 2) add the filter to the \texttt{allfilters.dat} file located in the
\texttt{data} directory. Follow the format: there must be a line starting with a
\# sign, followed by two columns, the first being the wavelength in angstroms,
the second being the total throughput.  The filter can be of any resolution, and
need not be properly normalized.  3) The user would be wise to add details of
the filter to the \texttt{FILTER\_LIST} file located in the \texttt{data}
directory, although this is not, strictly speaking, necessary for proper
functioning of the code.

We would appreciate it if the user would email us if they add a filter so that
we can include this filter in later releases (thereby saving others the trouble
of adding filters).

\vspace{0.4in}
\noindent{\bf 5.  Acknowledgements}
\vspace{0.1in}

FSPS has been in active development since 2009 and has been developed and tested
by a core group of users over the years including Martin White, Ben Johnson,
Nell Byler, Alexa Villaume, Joel Leja, and Demitri Muna.  Special thanks in
particular to Ben Johnson who has re-written entire routines and has extensively
tested the outputs.  I would also like to thank the many researchers who have
made publicly available their codes and/or models, which made it possible to
include many of the features in FSPS, including B. Aringer, Bruce Draine, Gary
Ferland, Karl Gordon, Zeljko Ivezic, Ariane Lancon, Maia Nenkova, Linda Smith,
the Padova, BaSTI, and Geneva stellar evolution groups, and the teams
responsible for the MILES and BaSeL spectral libraries including Bob Kurucz,
whose stellar atmosphere and spectral libraries underpine the BaSeL library and
the C3K spectral library under development.

The development of FSPS has been supported by Packard and Sloan foundation
fellowships, NASA grants NNX14AR86G and NNX15AK14G and NSF grant AST-1524161.


\clearpage

\begin{table}
\begin{center}
\caption{Lookup table of metallicity values depends on the
  isochrone set being used.}
\begin{tabular}{ccccccc}
\hline\hline
\texttt{zmet} & $Z$ [log($Z/Z_\odot$)] &  $Z$ [log($Z/Z_\odot$)] &
log($Z/Z_\odot$) & $Z$ [log($Z/Z_\odot$)] & $Z$ [log($Z/Z_\odot$)] & Z\\ &
Padova & BaSTI & MIST & Geneva & PARSEC & BPASS\\
\hline
1 &  0.0002  (-1.98) &  0.0003  (-1.82) & -2.50 & 0.0010 (-1.30) & 0.0001 (-2.18)& 0.0001\\
2 &  0.0003  (-1.80) &  0.0006  (-1.52) & -2.25 & 0.0040 (-0.70) & 0.0002 (-1.88)& 0.0010\\
3 &  0.0004  (-1.68) &  0.0010  (-1.30) & -2.00 & 0.0080 (-0.40) & 0.0005 (-1.48)& 0.0020\\
4 &  0.0005  (-1.58) &  0.0020  (-1.00) & -1.75 & 0.0200 (+0.00) & 0.0010 (-1.18)& 0.0030\\
5 &  0.0006  (-1.50) &  0.0040  (-0.70) & -1.50 & 0.0400 (+0.30) & 0.0020 (-0.88)& 0.0040\\
6 &  0.0008  (-1.38) &  0.0080  (-0.40) & -1.25 &                & 0.0040 (-0.58)& 0.0060\\
7 &  0.0010  (-1.28) &  0.0100  (-0.30) & -1.00 &                & 0.0060 (-0.40)& 0.0080\\
8 &  0.0012  (-1.20) &  0.0200  (+0.00) & -0.75 &                & 0.0080 (-0.28)& 0.0100\\
9 &  0.0016  (-1.07) &  0.0300  (+0.18) & -0.50 &                & 0.0100 (-0.18)& 0.0140\\
10 & 0.0020  (-0.98) &  0.0400  (+0.30) & -0.25 &                & 0.0140 (-0.04)& 0.0200\\
11 & 0.0025  (-0.89) &                  & +0.00 &                & 0.0170 (+0.05)& 0.0300\\
12 & 0.0031  (-0.79) &                  & +0.25 &                & 0.0200 (+0.12)& 0.0400\\
13 & 0.0039  (-0.69) &                  & +0.50 &                & 0.0300 (+0.30)& \\
14 & 0.0049  (-0.59) &  & & & 0.0400 (+0.42) &\\
15 & 0.0061  (-0.49) &  & & & 0.0600 (+0.60) &\\
16 & 0.0077  (-0.39) &  &\\
17 & 0.0096  (-0.30) &  &\\
18 & 0.0120  (-0.20) &  &\\
19 & 0.0150  (-0.10) &  &\\
20 & 0.0190  (+0.00) & & \\
21 & 0.0240  (+0.10) & & \\
22 & 0.0300  (+0.20) & & \\
\hline
\end{tabular}
\end{center}
\tablecomments{$Z_\odot=0.0190$ for Padova, $Z_\odot=0.0200$ for BaSTI
  and Geneva, $Z_\odot=0.0191$ for MIST, $Z_\odot=0.0152$ for PARSEC,
  and $Z_\odot=0.020$ for BPASS.}
\end{table}

% \clearpage

% \twocolumn

% \begin{table}
% \begin{center}
% \caption{Emission lines included in FSPS}
% \begin{tabular}{ccc}
% \hline\hline
% 923.148 & Ly 923 & \texttt{H  1 923.156A}\\
% 926.249 & Ly 926 & \texttt{H  1 926.231A}\\
% 930.751 & Ly 930 & \texttt{H  1 930.754A}\\
% 937.814 & Ly 937 & \texttt{H  1 937.809A}\\
% 949.742 & Ly-$\delta$ 949 & \texttt{H  1 949.749A}\\
% 972.517 & Ly-$\gamma$ 972 & \texttt{H  1 972.543A}\\
% 1025.728 & Ly-$\beta$ 1025 & \texttt{H  1 1025.73A}\\
% 1215.6701 & Ly-$\alpha$ 1216 & \texttt{H  1 1215.68A}\\
% 1640.42 & He{\sc\,ii} 1640 & \texttt{He 2 1640.00A}\\
% 1661.241 & O{\sc\,iii}] 1661 & \texttt{O  3 1661.00A}\\
% 1666.15 & O{\sc\,iii}] 1666 & \texttt{O  3 1666.00A}\\
% 1812.205 & [Ne{\sc\,iii}] 1815 & \texttt{Ne 3 1815.00A}\\
% 1854.716 & [Al{\sc\,iii}] 1855 & \texttt{Al 3 1855.00A}\\
% 1862.7895 & [Al{\sc\,iii}] 1863 & \texttt{Al 3 1863.00A}\\
% 1906.68 & [C{\sc\,iii}]  & \texttt{C  3 1907.00A}\\
% 1908.73 & [C{\sc\,iii}]  & \texttt{C  3 1910.00A}\\
% 2142.3 & N{\sc\,ii}] 2141 & \texttt{N  2 2141.00A}\\
% 2321.664 & [O{\sc\,iii}] 2321 & \texttt{O  3 2321.00A}\\
% 2324.21 & C{\sc\,ii}] 2326 & \texttt{C  2 2324.00A}\\
% 2325.4 & C{\sc\,ii}] 2326 & \texttt{C  2 2325.00A}\\
% 2326.11 & C{\sc\,ii}] 2326 & \texttt{C  2 2327.00A}\\
% 2327.64 & C{\sc\,ii}] 2326 & \texttt{C  2 2328.00A}\\
% 2328.83 & C{\sc\,ii}] 2326 & \texttt{C  2 2329.00A}\\
% 2471.088 & [O{\sc\,ii}] 2471 & \texttt{O II 2471.00A}\\
% 2661.146 & [Al{\sc\,ii}] 2660 & \texttt{Al 2 2660.00A}\\
% 2669.951 & [Al{\sc\,ii}] 2670 & \texttt{Al 2 2670.00A}\\
% 2796.352 & Mg{\sc\,ii} 2800 & \texttt{Mg 2 2795.53A}\\
% 2803.53 & Mg{\sc\,ii} 2800 & \texttt{Mg 2 2802.71A}\\
% 3109.98 & [Ar{\sc\,iii}] 3110 & \texttt{Ar 3 3109.00A}\\
% 3343.5 & [Ne{\sc\,iii}] 3343 & \texttt{Ne 3 3343.00A}\\
% 3722.75 & [S{\sc\,iii}] 3723 & \texttt{S  3 3722.00A}\\
% 3727.1 & [O{\sc\,ii}] 3726 & \texttt{O II 3726.00A}\\
% 3729.86 & [O{\sc\,ii}] 3729 & \texttt{O II 3729.00A}\\
% 3798.987 & H 3798 & \texttt{H  1 3797.92A}\\
% 3836.485 & H 3835 & \texttt{H  1 3835.40A}\\
% 3869.86 & [Ne{\sc\,iii}] 3870 & \texttt{Ne 3 3869.00A}\\
% 3889.75 & He{\sc\,i} 3889 & \texttt{He 1 3888.63A}\\
% 3890.166 & H 3889 & \texttt{H  1 3889.07A}\\
% 3968.59 & [Ne{\sc\,iii}] 3968 & \texttt{Ne 3 3968.00A}\\
% 3971.198 & H 3970 & \texttt{H  1 3970.09A}\\
% 4069.75 & [S{\sc\,ii}] 4070 & \texttt{S II 4070.00A}\\
% 4077.5 & [S{\sc\,ii}] 4078 & \texttt{S II 4078.00A}\\
% 4102.892 & H-$\delta$ 4102 & \texttt{H  1 4101.76A}\\
% 4341.692 & H-$\gamma$ 4340 & \texttt{H  1 4340.49A}\\
% 4364.435 & [O{\sc\,iii}] 4364 & \texttt{TOTL 4363.00A}\\
% 4472.735 & He{\sc\,i} 4472 & \texttt{He 1 4471.47A}\\
% 4622.864 & [C{\sc\,i}] 4621 & \texttt{C  1 4621.00A}\\
% 4725.47 & [Ne{{\sc\,iv}}] 4720 & \texttt{Ne 4 4720.00A}\\
% 4862.71 & H-$\beta$ 4861 & \texttt{H  1 4861.36A}\\
% 4960.295 & [O{\sc\,iii}] 4960 & \texttt{O  3 4959.00A}\\
% 5008.24 & [O{\sc\,iii}] 5007 & \texttt{O  3 5007.00A}\\
% 5193.27 & [Ar{\sc\,iii}] 5193 & \texttt{Ar 3 5192.00A}\\
% 5201.705 & [N{\sc\,i}] 5200 & \texttt{N  1 5200.00A} \\
% \hline
% \end{tabular}
% \end{center}
% \end{table}

% \begin{table}
% \begin{center}
% \caption{Emission lines included in FSPS (cont.)}
% \begin{tabular}{ccc}
% \hline\hline
% 5519.242 & [Cl{\sc\,iii}] 5518 & \texttt{Cl 3 5518.00A}\\
% 5539.411 & [Cl{\sc\,iii}] 5538 & \texttt{Cl 3 5538.00A}\\
% 5578.89 & [O{\sc\,i}] 5578 & \texttt{O  1 5577.00A}\\
% 5756.19 & [N{\sc\,ii}] 5756 & \texttt{N  2 5755.00A}\\
% 5877.249 & He{\sc\,i} 5877 & \texttt{He 1 5875.61A}\\
% 6302.046 & [O{\sc\,i}] 6302 & \texttt{O  1 6300.00A}\\
% 6313.81 & [S{\sc\,iii}] 6314 & \texttt{S  3 6312.00A}\\
% 6365.535 & [O{\sc\,i}] 6365 & \texttt{O  1 6363.00A}\\
% 6549.86 & [N{\sc\,ii}] 6549 & \texttt{N  2 6548.00A}\\
% 6564.6 & H-$\alpha$ 6563 & \texttt{H  1 6562.85A}\\
% 6585.27 & [N{\sc\,ii}] 6585 & \texttt{N  2 6584.00A}\\
% 6679.995 & He{\sc\,i} 6680 & \texttt{He 1 6678.15A}\\
% 6718.294 & [S{\sc\,ii}] 6717 & \texttt{S II 6716.00A}\\
% 6732.673 & [S{\sc\,ii}] 6732 & \texttt{S II 6731.00A}\\
% 7067.138 & He{\sc\,i} 7065 & \texttt{He 1 7065.18A}\\
% 7137.77 & [Ar{\sc\,iii}] 7138 & \texttt{Ar 3 7135.00A}\\
% 7321.94 & [O{\sc\,ii}] 7323 & \texttt{O II 7323.00A}\\
% 7332.21 & [O{\sc\,ii}] 7332 & \texttt{O II 7332.00A}\\
% 7334.17 & [Ar{{\sc\,iv}}] 7330 & \texttt{Ar 4 7331.00A}\\
% 7753.19 & [Ar{\sc\,iii}] 7753 & \texttt{Ar 3 7751.00A}\\
% 8581.06 & [Cl{\sc\,ii}] 8579 & \texttt{Cl 2 8579.00A}\\
% 8729.53 & [C{\sc\,i}] 8727 & \texttt{C  1 8727.00A}\\
% 9017.8 & Pa 9015 & \texttt{H  1 9014.92A}\\
% 9071.1 & [S{\sc\,iii}] 9071 & \texttt{S  3 9069.00A}\\
% 9126.1 & [Cl{\sc\,ii}] 9124 & \texttt{Cl 2 9124.00A}\\
% 9232.2 & Pa 9229 & \texttt{H  1 9229.03A}\\
% 9533.2 & [S{\sc\,iii}] 9533 & \texttt{S  3 9532.00A}\\
% 9548.8 & Pa 9546 & \texttt{H  1 9545.99A}\\
% 9852.96 & [C{\sc\,i}] 9850 & \texttt{TOTL 9850.00A}\\
% 10052.6 & Pa-$\delta$ 10050 & \texttt{H  1 1.00494m}\\
% 10323.32 & [S{\sc\,ii}] 10331 & \texttt{S  2 1.03300m}\\
% 10832.057 & He{\sc\,i} 10829 & \texttt{He 1 1.08299m}\\
% 10833.306 & He{\sc\,i} 10833 & \texttt{He 1 1.08303m}\\
% 10941.17 & Pa-$\gamma$ 10939 & \texttt{H  1 1.09381m}\\
% 12570.21 & [Fe{\sc\,ii}] 1.26$\mu\mathrm{m}$ & \texttt{Fe 2 1.25668m}\\
% 12821.578 & Pa-$\beta$ 12819 & \texttt{H  1 1.28181m}\\
% 17366.885 & Br 17363 & \texttt{H  1 1.73621m}\\
% 18179.2 & Br 18175 & \texttt{H  1 1.81741m}\\
% 18756.4 & Pa-$\alpha$ 18752 & \texttt{H  1 1.87511m}\\
% 19450.89 & Br-$\delta$ 19447 & \texttt{H  1 1.94456m}\\
% 21661.178 & Br-$\gamma$ 21657 & \texttt{H  1 2.16553m}\\
% 26258.71 & Br-$\beta$ 26254 & \texttt{H  1 2.62515m}\\
% 30392.02 & Pf 30386 & \texttt{H  1 3.03837m}\\
% 32969.8 & Pf-$\delta$ 32964 & \texttt{H  1 3.29609m}\\
% 37405.76 & Pf-$\gamma$ 37398 & \texttt{H  1 3.73953m}\\
% 40522.79 & Br-$\alpha$ 40515 & \texttt{H  1 4.05116m}\\
% 46537.8 & Pf-$\beta$ 46529 & \texttt{H  1 4.65250m}\\
% 51286.5 & Hu-$\delta$ 51277 & \texttt{H  1 5.12725m}\\
% 59082.2 & Hu-$\gamma$ 59071 & \texttt{H  1 5.90659m}\\
% 69852.74 & [Ar{\sc\,ii}] 7$\mu\mathrm{m}$ & \texttt{Ar 2 6.98000m}\\
% 74599.0 & Pf-$\alpha$ 74585 & \texttt{H  1 7.45781m}\\
% 75024.4 & Hu-$\beta$ 75011 & \texttt{H  1 7.50043m}\\
% 89913.8 & [Ar{\sc\,iii}] 9$\mu\mathrm{m}$ & \texttt{Ar 3 9.00000m}\\
% \hline
% \end{tabular}
% \end{center}
% \end{table}

% \begin{table}
% \begin{center}
% \caption{Emission lines included in FSPS (cont.)}
% \begin{tabular}{ccc}
% \hline\hline
% 105105.0 & [S{{\sc\,iv}}] 10.5$\mu\mathrm{m}$ & \texttt{S  4 10.5100m}\\
% 123719.12 & Hu-$\alpha$ 12.4$\mu\mathrm{m}$ & \texttt{H  1 12.3685m}\\
% 128135.48 & [Ne{\sc\,ii}] 12.8$\mu\mathrm{m}$ & \texttt{Ne 2 12.8100m}\\
% 143678.0 & [Cl{\sc\,ii}] 14.4$\mu\mathrm{m}$ & \texttt{Cl 2 14.4000m}\\
% 155551.0 & [Ne{\sc\,iii}] 15.5$\mu\mathrm{m}$ & \texttt{Ne 3 15.5500m}\\
% 187130.0 & [S{\sc\,iii}] 18.7$\mu\mathrm{m}$ & \texttt{S  3 18.6700m}\\
% 218302.0 & [Ar{\sc\,iii}] 22$\mu\mathrm{m}$ & \texttt{Ar 3 21.8300m}\\
% 328709.0 & [P{\sc\,ii}] 32$\mu\mathrm{m}$ & \texttt{P  2 32.8700m}\\
% 334800.0 & [S{\sc\,iii}] 33.5$\mu\mathrm{m}$ & \texttt{S  3 33.4700m}\\
% 348140.0 & [Si{\sc\,ii}] 35$\mu\mathrm{m}$ & \texttt{Si 2 34.8140m}\\
% 360135.0 & [Ne{\sc\,iii}] 36$\mu\mathrm{m}$ & \texttt{Ne 3 36.0140m}\\
% 518145.0 & [O{\sc\,iii}] 52$\mu\mathrm{m}$ & \texttt{O  3 51.8000m}\\
% 573300.0 & [N{\sc\,iii}] 57$\mu\mathrm{m}$ & \texttt{N  3 57.2100m}\\
% 606420.0 & [P{\sc\,ii}] 60$\mu\mathrm{m}$ & \texttt{P  2 60.6400m}\\
% 631852.0 & [O{\sc\,i}] 63$\mu\mathrm{m}$ & \texttt{O  1 63.1700m}\\
% 883564.0 & [O{\sc\,iii}] 88$\mu\mathrm{m}$ & \texttt{O  3 88.3300m}\\
% 1218000.0 & [N{\sc\,ii}] 122$\mu\mathrm{m}$ & \texttt{N  2 121.700m}\\
% 1455350.0 & [O{\sc\,i}] 145$\mu\mathrm{m}$ & \texttt{O  1 145.530m}\\
% 1576429.62 & [C{\sc\,ii}] 157.7$\mu\mathrm{m}$ & \texttt{C  2 157.600m}\\
% 2053000.0 & [N{\sc\,ii}] 205$\mu\mathrm{m}$ & \texttt{N  2 205.400m}\\
% 3703700.0 & [C{\sc\,i}] 369$\mu\mathrm{m}$ & \texttt{C  1 369.700m}\\
% 6097000.0 & [C{\sc\,i}] 610$\mu\mathrm{m}$ & \texttt{C  1 609.200m} \\
% \hline
% \end{tabular}
% \end{center}
% \end{table}



\end{document}











































